\chapter{Cross-Sectional and Cross-Asset Features}\label{ch:rs:cross-sectional-and-cross-asset-features}

A firm sells much of its output to one customer. The customer reports bad news, and the firm's prospects have just
worsened with it; its stock should fall the same day. Cohen and Frazzini, using a dataset of firms' principal customers, found
that stock prices do not incorporate news involving related firms promptly, generating predictable subsequent price moves: a
long--short strategy built on the effect earned monthly alphas of over 150 basis points. A feature built from other securities
asks one question the security's own history cannot: who moves first, and does the other one follow? This chapter measures it
three ways, at three clocks: US size portfolios over sixty years of weeks; a synthetic market with planted customer--supplier
links, over months; and two simulated instruments a planted fraction of a second apart, over ticks. One lesson recurs at every
clock: the leader's past must predict the follower beyond the follower's own past.

\section{Peer and sector residuals}

\begin{definition}[Peer group, peer-relative return]\label{def:rs:cross-sectional-and-cross-asset-features:peer}
A \emph{peer group}\index{peer group} of a security is a set of securities chosen to share its common drivers: its industry,
its sector, or its nearest neighbours by return correlation or by business description. Its \emph{peer-relative
return}\index{peer-relative return} is its return minus the mean return of the other members of its group over the same
period.
\end{definition}

Subtracting peers is the cheapest residualisation, a special case of the \omterm{def:rs:anatomy-of-a-predictor:residual}{residual return} of chapter~6 with the peer mean as
the only factor, exposure one. The word \emph{other} matters.

\begin{proposition}[A security is not its own peer]\label{prop:rs:cross-sectional-and-cross-asset-features:loo}
In a group of $n$ securities, the return minus the group mean \emph{including} the security equals $(n-1)/n$ times the return
minus the mean of the other $n-1$.
\end{proposition}

\begin{proof}
With $m_{-i}$ the mean of the others, the full mean is $(x_i + (n-1)m_{-i})/n$, and $x_i - (x_i + (n-1)m_{-i})/n =
\frac{n-1}{n}(x_i - m_{-i})$.
\end{proof}

The shrinkage is harmless for a rank but not for a level: in a group of four, a \omterm{def:rs:cross-sectional-and-cross-asset-features:peer}{peer-relative return} that includes the
security is a quarter too small, and the error differs across groups of different sizes, so a threshold or a pooled regression
reads groups unequally. \texttt{firm.leadlag.peer\_relative} leaves the security out.

Peers matter for cross-sectional features built from \emph{other} securities' returns. When a customer and its supplier share a
factor (their industry, the market), the factor part of the customer's move reaches the supplier the same day and is priced;
only the customer's own news can arrive late. The peer-relative customer return isolates that news. How much it helps depends on
how much of the customer's move is common; the next section measures a case.

\section{Lead--lag at daily horizons}

\begin{definition}[Lead--lag relationship, economic link]\label{def:rs:cross-sectional-and-cross-asset-features:leadlag}
Two securities have a \emph{lead--lag relationship}\index{lead--lag relationship} when the past returns of one (the leader)
predict the future returns of the other (the follower) beyond what the follower's own past predicts. An \emph{economic
link}\index{economic link} is a documented business relationship between two firms (customer and supplier, competitors,
partners, a parent and its subsidiary) that gives a reason to expect one.
\end{definition}

The definition asks for predictability beyond the follower's own past, the Granger causality of Book~4, chapter~20, and the
clause is not a formality. Lo and MacKinlay (1990) showed important lead--lag relations across US securities;
Hou (2007) located them within industries, big firms leading small ones. The data library of Kenneth French has daily returns of
US size quintiles back to 1926, from which \texttt{rs\_fetch\_size.py} stores only derived statistics (the library carries a
copyright notice and no licence). Over weeks, the correlation of the small quintile's return with the large quintile's previous
week was 0.31 from 1962 to 1987, 0.23 from 1988 to 2006 and 0.06 from 2007 to mid-2026; the reverse correlation, small leading
large, was 0.06, $-0.09$ and $-0.05$ (\cref{fig:rs:cross-sectional-and-cross-asset-features:size}).

\begin{omfigure}
\begin{tikzpicture}
\begin{axis}[width=11cm,height=5.4cm,ybar,bar width=10pt,ylabel={correlation},tick label style={font=\small},
  label style={font=\small},xtick=data,xticklabels={1962--1987,1988--2006,2007--2026},ymin=-0.15,ymax=0.45,
  enlarge x limits=0.25,grid=major,grid style={gray!25},legend columns=3,
  legend style={font=\footnotesize,at={(0.5,-0.2)},anchor=north,draw=none,fill=none,
  /tikz/every even column/.append style={column sep=8pt}},table/col sep=comma]
\addplot[fill=omDef!55,draw=omDef] table[x=k,y=small_on_big]{figdata/research/10-cross-sectional-and-cross-asset-features/size.csv};
\addplot[fill=omThm!45,draw=omThm] table[x=k,y=big_on_small]{figdata/research/10-cross-sectional-and-cross-asset-features/size.csv};
\addplot[fill=gray!40,draw=gray] table[x=k,y=own_small]{figdata/research/10-cross-sectional-and-cross-asset-features/size.csv};
\legend{small on big's last week,big on small's last week,small on its own last week}
\end{axis}
\end{tikzpicture}

\omcaption{Weekly lead--lag between the smallest and the largest size quintiles of US stocks (equal-weighted), in three periods:
the correlation of each quintile's return with the other's, and with its own, previous week. Derived from the Kenneth R. French
Data Library (daily size portfolios, 202607 CRSP file).}\label{fig:rs:cross-sectional-and-cross-asset-features:size}
\end{omfigure}

The grey \omterm{def:rs:market-data-for-research:bar}{bars} are the catch. The small quintile's return was as correlated with its own previous week (0.37) as with the large
quintile's (0.31), and a regression of the small quintile's week on both previous weeks gives the large quintile a $t$ statistic
of 1.9 from 1962 to 1987, $-0.7$ from 1988 to 2006 and 1.1 since: over weeks, most of the famous lead is the small stocks' own
autocorrelation, the kind of short-horizon autocorrelation that Boudoukh, Richardson and Whitelaw (1994) argued had been overstated
in the literature and traced most likely to institutional factors. Over days the lead is sharper: from 1962 to 1987 the large quintile's previous
day has a $t$ statistic of 13.4 beyond the small quintile's own; from 2007 on it is $-4.4$. Either way, it has faded.

\subsection{Economic links, planted}

A synthetic market can hold links whose truth is known. \texttt{firm.synthmkt} with \texttt{link\_share = 0.3} makes 30\% of
listings suppliers of one customer drawn at random, and adds to each supplier's return 10\% of its customer's specific shocks,
spread evenly over the next 21 trading days: an attention delay in the spirit of Cohen and Frazzini. Around the customer's earnings
announcements (\cref{fig:rs:cross-sectional-and-cross-asset-features:event}) the customer moves 4.5\% per unit of surprise on the
day, and the supplier drifts to 0.39\% over the next 21 days.

\begin{omfigure}
\begin{tikzpicture}
\begin{axis}[width=11cm,height=5.4cm,xlabel={trading days from the customer's announcement},
  ylabel={cumulative return per unit surprise (\%)},tick label style={font=\small},label style={font=\small},xmin=-5,xmax=40,
  ymin=-0.2,ymax=0.6,grid=major,grid style={gray!25},legend pos=south east,legend style={font=\footnotesize},
  table/col sep=comma]
\addplot[omThm,thick,dashed] table[x=day,y=customer_tenth]{figdata/research/10-cross-sectional-and-cross-asset-features/event.csv};
\addplot[omDef,very thick] table[x=day,y=supplier]{figdata/research/10-cross-sectional-and-cross-asset-features/event.csv};
\legend{customer ($\times 0.1$),supplier}
\end{axis}
\end{tikzpicture}

\omcaption{Market-adjusted cumulative returns around customers' earnings announcements, as slopes on the true surprise across
11\,100 customer events and 9\,654 supplier events: the customer's move (divided by ten, the planted response) and its suppliers'.
The supplier gets there over a month. Data: \texttt{firm.synthmkt}, \texttt{link\_share = 0.3}, seed 1.}\label{fig:rs:cross-sectional-and-cross-asset-features:event}
\end{omfigure}

The feature follows: at each month-end, each supplier's customer return over the past month. Its rank IC with the supplier's next
month is 0.035 over 107 months ($t = 5.6$), and a long--short of the top and bottom quintiles earns 0.90\% a month with an
annualised Sharpe ratio of 1.9. The truth (the planted link part of each supplier's expected return over the coming month) has an
IC of 0.074: the customer's last month is a noisy proxy for the part of its news still to diffuse, and half the information is
lost in the proxy. The peer-relative customer return does barely better (0.036), because the simulator plants the link on specific
shocks only and its industry factors are small; in data where customers and suppliers share large common moves the difference is
the reason to use it.

\subsection{Mining links without the list}

Cohen and Frazzini had the list of customers. Without it, a researcher can rank every ordered pair of securities by the correlation
of one's return with the other's previous return and hope the true links rise to the top. On the synthetic market, the 698 names
listed for all ten years form 486\,506 ordered pairs, of which 133 are planted links:

\begin{center}
\small
\begin{tabular}{@{}lccc@{}}
\toprule
 & links in the top 100 & in the top 1\,000 & in the top 10\,000\\
\midrule
daily returns, one-day lag & 0 & 0 & 3\\
monthly returns, one-month lag & 0 & 2 & 8\\
at random (expected) & 0.03 & 0.27 & 2.7\\
\bottomrule
\end{tabular}
\end{center}

A true pair's lagged monthly correlation averages 0.052 while the correlations of all pairs have a standard deviation of 0.097: a
link sits half a standard deviation above the noise, and among half a million pairs the noise wins. The customer list finds all 133
links by construction. Mining a lead--lag network is the multiple-testing problem of chapter~20 in its purest form; the \omterm{def:rs:cross-sectional-and-cross-asset-features:leadlag}{economic
link} is the prior that makes the search small enough to succeed.

\section{Futures-to-cash and cross-venue}

\begin{definition}[Futures-to-cash lead, lead--lag estimator]\label{def:rs:cross-sectional-and-cross-asset-features:estimator}
The \emph{futures-to-cash lead}\index{futures-to-cash lead} is the tendency of an index future's price to move before the prices
of the index's constituents and of the funds that track it. A \emph{lead--lag estimator}\index{lead--lag estimator} estimates the
time shift $\vartheta$ at which two price series are most related: here the lag maximising the Hayashi--Yoshida cross-covariance of
one series with the other's clock shifted by $\vartheta$, computed from all observations without synchronising them (Hoffmann,
Rosenbaum and Yoshida, 2013).
\end{definition}

At high frequency, Huth and Abergel (2014) observed strongly asymmetric cross-correlation functions, especially between futures
and stocks; the most liquid assets tend to lead. They reached 60\% accuracy forecasting the next mid-quote move of the lagger from
the leader's past, significantly better than from the lagger's own; a naive strategy with market orders could not profit from it
because of the bid--ask spread. \texttt{firm.tape.simulate\_pair} builds the case with a known answer: two instruments on one
efficient price, the second's price delayed by a planted latency, each with its own liquidity and order flow. With the efficient
price jumping about once a second, each mid changes 0.65 times a second.

\begin{omfigure}
\begin{tikzpicture}
\begin{axis}[width=11cm,height=5.2cm,xlabel={lag $\vartheta$ (seconds)},ylabel={cross-correlation},tick label style={font=\small},
  label style={font=\small},xmin=-3,xmax=3,ymin=0.1,ymax=0.25,grid=major,grid style={gray!25},legend columns=2,
  legend style={font=\footnotesize,at={(0.5,-0.3)},anchor=north,draw=none,fill=none,
  /tikz/every even column/.append style={column sep=8pt}},table/col sep=comma]
\addplot[omDef,very thick] table[x=lag,y=fast]{figdata/research/10-cross-sectional-and-cross-asset-features/ccf.csv};
\addplot[gray,thick,dashed] table[x=lag,y=slow]{figdata/research/10-cross-sectional-and-cross-asset-features/ccf.csv};
\draw[omThm,thick] (axis cs:0.5,0.1) -- (axis cs:0.5,0.25);
\node[font=\footnotesize,omThm,anchor=west] at (axis cs:0.55,0.12) {planted 0.5 s};
\legend{fast pair (0.65 changes/s),slow pair (0.16 changes/s)}
\end{axis}
\end{tikzpicture}

\omcaption{Hayashi--Yoshida cross-correlation of two simulated instruments' mid-prices, the second's clock shifted by $\vartheta$,
for a planted latency of 0.5 seconds; normalised by one-minute realised variances. The fast pair's function is flat-topped and
centred on the latency; the slow pair's price changes too rarely to show it. Data: \texttt{firm.tape}, one hour, seed
10.}\label{fig:rs:cross-sectional-and-cross-asset-features:ccf}
\end{omfigure}

The cross-correlation is flat near its top (\cref{fig:rs:cross-sectional-and-cross-asset-features:ccf}): each mid follows its
efficient price with a spread of delays, so the peak is smeared over a second either way, and its argmax lands at 0.15 seconds
for a planted 0.5. The smear is symmetric when both instruments react alike, so the centre of symmetry is a better estimate than
the peak: \texttt{lead\_symmetric} takes the lag about which the function is most symmetric over one second each side, and
finds 0.5. Over five planted latencies and three seeds (\cref{fig:rs:cross-sectional-and-cross-asset-features:recovery}), the
centre misses by 0.053 seconds on average and 0.2 at worst, the argmax by 0.19 on average and 0.5 at worst. On the slow pair,
whose mids change every six seconds, both fail: a lead shorter than the time between price changes is invisible to any estimator
built on price changes.

\begin{omfigure}
\begin{tikzpicture}
\begin{axis}[width=9cm,height=5.6cm,xlabel={planted latency (seconds)},ylabel={estimated lead (seconds)},
  tick label style={font=\small},label style={font=\small},xmin=-0.2,xmax=2.2,ymin=-1,ymax=2.6,grid=major,grid style={gray!25},
  legend pos=north west,legend style={font=\footnotesize},table/col sep=comma]
\addplot[black!50,domain=-0.2:2.2,samples=2,forget plot] {x};
\addplot[only marks,mark=*,omDef,mark size=2.2pt] table[x=planted,y=fast]{figdata/research/10-cross-sectional-and-cross-asset-features/recovery.csv};
\addplot[only marks,mark=x,gray,mark size=3pt,thick] table[x=planted,y=slow]{figdata/research/10-cross-sectional-and-cross-asset-features/recovery.csv};
\legend{fast pair,slow pair}
\end{axis}
\end{tikzpicture}

\omcaption{Estimated lead (symmetric centre of the Hayashi--Yoshida cross-correlation) against the planted latency, three seeds each,
with the diagonal. Data: \texttt{firm.tape}, one hour per run.}\label{fig:rs:cross-sectional-and-cross-asset-features:recovery}
\end{omfigure}

A feature needs sampled returns. At each sampling interval, the table gives the correlation of the follower's return with the
leader's previous return, the reverse, and the follower's correlation with its own previous return, for the fast pair with the
planted 0.5 seconds:

\begin{center}
\small
\begin{tabular}{@{}lccccccc@{}}
\toprule
interval (s) & 0.1 & 0.5 & 1 & 2 & 5 & 10 & 30\\
\midrule
follower on leader's last & 0.12 & 0.39 & 0.55 & 0.59 & 0.48 & 0.26 & $-0.02$\\
leader on follower's last & 0.11 & 0.33 & 0.43 & 0.45 & 0.31 & 0.16 & $-0.04$\\
follower on its own last & 0.13 & 0.31 & 0.40 & 0.44 & 0.35 & 0.22 & $-0.03$\\
\bottomrule
\end{tabular}
\end{center}

Every row is large, and that is the simulator: its mids reach the efficient price in steps over a few seconds, so all mid-price
returns trend at short intervals (real mid-quote returns do not, to anything like this degree). The planted lead is the
\emph{asymmetry}, the first row minus the second: 0.13 at two seconds, 0.17 at five, 0.10 at ten, nothing at thirty seconds;
and the first row must also beat the third, which at two seconds it does by 0.14. The lead's footprint is a narrow band of sampling intervals, a few times the latency; coarser sampling averages it away,
and finer sampling sees too few price changes. The same test that dismantled the weekly size lead applies: the leader must beat
the follower's own past.

\section{Cross-asset signals}

\begin{definition}[Cross-asset feature]\label{def:rs:cross-sectional-and-cross-asset-features:crossasset}
A \emph{cross-asset feature}\index{cross-asset feature} predicts a security from the prices of a different asset class or market:
an index future for its constituents, a currency for its exporters, a commodity for its producers, rates for equities, one market's
close for another market's open.
\end{definition}

Each needs a mechanism, and the mechanism decides the clock. An arbitrage (a future and its basket, a fund and its holdings, the
fair value of Book~1, chapter~21) holds the prices together and leaves a lead of milliseconds to seconds, fought over by the fastest (Book~11). An economic exposure (a producer's revenue in a commodity, an exporter's in a currency) leaves room for
slow diffusion over days or months, the timescale of the customer link. Two traps are specific to them. The first is clocks: a
market that closes later carries information the earlier one could not have had, so a Tokyo close regressed on a New York close of
the same date looks predictive and is a leak; each value must be stamped with the time it became known (chapter~3). The second is
the follower's own past, the lesson of both sections above.

\section{Predictor cards}

\begin{predictorcard}{Customer momentum}\label{pred:rs:cross-sectional-and-cross-asset-features:customer}
\sfield{Definition} For each supplier, its principal customer's return over the last month, peer-relative; ranked across suppliers.
\sfield{Inputs and timestamps} A customer--supplier list as known at the time (disclosures are filed with a lag: chapter~3); the
customers' returns to the month-end close.
\sfield{Rationale} Attention-constrained investors do not incorporate news about related firms promptly (Cohen and Frazzini).
\sfield{Horizon and half-life} One month; on the synthetic market with a 21-day planted diffusion, IC 0.035 (truth 0.074).
\sfield{Normalisation} Rank across suppliers; neutralise the supplier's own industry and momentum.
\sfield{Failure modes} Stale or look-ahead link data; links that are shared factor exposure rather than news; crowding once the
effect is published (chapter~28).
\sfield{Sources} Cohen and Frazzini (2008); Menzly and Ozbas (2010); \texttt{rs\_leadlag.customer\_momentum}.
\end{predictorcard}

\begin{predictorcard}{Leader's last return}\label{pred:rs:cross-sectional-and-cross-asset-features:leader}
\sfield{Definition} The leader's mid-price return over the last interval $\Delta$, as a \omterm{def:rs:anatomy-of-a-predictor:predictor}{predictor} of the follower's return over the next.
\sfield{Inputs and timestamps} Both instruments' quotes, stamped on one clock at the point of receipt (the latency between venues is the
signal and must not be mistaken for a clock difference).
\sfield{Rationale} Price discovery happens first where trading is cheapest and most liquid (Huth and Abergel).
\sfield{Horizon and half-life} A few times the lead; on the simulated pair, the excess over the reverse direction is largest at two to
five seconds and gone at thirty.
\sfield{Normalisation} By the follower's volatility over the interval.
\sfield{Failure modes} The follower's own autocorrelation mistaken for a lead; a lead too short for the price-change rate; the spread
(the signal may not pay for crossing it).
\sfield{Sources} Hoffmann, Rosenbaum and Yoshida (2013); Huth and Abergel (2014); \texttt{rs\_leadlag.interval\_table}.
\end{predictorcard}

\section{Tutorial: who moves first}\label{tut:rs:cross-sectional-and-cross-asset-features:tut}

\begin{tutorial}
\textbf{Goal.} Estimate the lead of one simulated instrument over another from their quotes; recover planted customer--supplier links
in a synthetic market and see what mining finds without them. \textbf{End state:}
\cref{fig:rs:cross-sectional-and-cross-asset-features:ccf,fig:rs:cross-sectional-and-cross-asset-features:recovery,fig:rs:cross-sectional-and-cross-asset-features:event};
the tables of this chapter.
\begin{tutsteps}
\item \textbf{The cross-correlation at lags.} Book 4's Hayashi--Yoshida sum with the second clock shifted, normalised by one-minute
realised variances (tick-level variances are distorted by the price's steps).
\omcode{code/firm/leadlag/firm_leadlag.py}{31}{42}{Hayashi--Yoshida cross-correlation at lags.}\label{lst:rs:cross-sectional-and-cross-asset-features:ccf}
\item \textbf{The symmetric centre}, among the lags where the function is at least half its maximum.
\omcode{code/firm/leadlag/firm_leadlag.py}{49}{56}{The lag of best symmetry.}\label{lst:rs:cross-sectional-and-cross-asset-features:centre}
\item \textbf{Peers without the security itself} (\cref{prop:rs:cross-sectional-and-cross-asset-features:loo}).
\omcode{code/firm/leadlag/firm_leadlag.py}{71}{83}{Leave-one-out peer-relative returns.}\label{lst:rs:cross-sectional-and-cross-asset-features:peer}
\item \textbf{Run} \texttt{lead\_recovery()}, \texttt{interval\_table()}, \texttt{customer\_momentum()}, \texttt{event\_study()},
\texttt{mining()} and \texttt{fig\_leadlag.py}; \texttt{rs\_fetch\_size.py} once, for the size statistics.
\end{tutsteps}
\textbf{What to change next.} Give the two instruments different reaction speeds (a follower whose market makers are slower) and watch
the ccf lose its symmetry; plant links within industries and compare raw and peer-relative customer returns.
\end{tutorial}

\section{Build: the lead--lag toolkit}\label{bld:rs:cross-sectional-and-cross-asset-features:leadlag}

\begin{build}
\bfield{Purpose} Lead--lag measurement on tick data and on daily panels, and the cross-sectional plumbing (peers, links) that features
of Books 8 and 9 (index arbitrage, pairs, customer momentum) are built from.
\bfield{Interface} \texttt{hy\_ccf(t1, x1, t2, x2, lags, norm\_step)}, \texttt{lead\_estimate}, \texttt{lead\_symmetric}, \texttt{lead\_lag\_ratio},
\texttt{directional\_corr}, \texttt{peer\_relative(ret, group)}, \texttt{linked\_feature(values, link)}, \texttt{lagged\_corr\_matrix},
\texttt{top\_pairs}; \texttt{firm.synthmkt} gains \texttt{link\_share}, \texttt{link\_beta}, \texttt{link\_days} and \texttt{Panel.customer}
(off by default, so earlier results are unchanged).
\bfield{Rules} A positive lag means the second series follows; peers never include the security; lagged correlations are for ranking
candidates, never for accepting them.
\bfield{Acceptance tests} \texttt{code/firm/leadlag/tests/}: a Brownian path observed at random times by two series with a known delay
(0, 0.7 and 2 seconds) is recovered within 0.21 seconds by both estimators; the directional correlation is one-sided; a noisy flat-topped
function is centred correctly; leave-one-out peers by hand; a planted one-day follower is the top mined pair.
\bfield{Stretch} Lead--lag by time of day; a multivariate version (one leader, many followers); link lists as \omterm{def:rs:point-in-time-data-and-the-biases:lookahead}{point-in-time data}
(\texttt{firm.pit}).
\end{build}

\begin{omsources}
\item L.~Cohen and A.~Frazzini, ``Economic links and predictable returns'', \emph{Journal of Finance} 63(4), 2008.
\item L.~Menzly and O.~Ozbas, ``Market segmentation and cross-predictability of returns'', \emph{Journal of Finance} 65(4), 2010.
\item A.~W.~Lo and A.~C.~MacKinlay, ``When are contrarian profits due to stock market overreaction?'', \emph{Review of Financial Studies}
3(2), 1990.
\item K.~Hou, ``Industry information diffusion and the lead-lag effect in stock returns'', \emph{Review of Financial Studies} 20(4), 2007.
\item J.~Boudoukh, M.~P.~Richardson and R.~F.~Whitelaw, ``A tale of three schools: insights on autocorrelations of short-horizon stock
returns'', \emph{Review of Financial Studies} 7(3), 1994.
\item M.~Hoffmann, M.~Rosenbaum and N.~Yoshida, ``Estimation of the lead-lag parameter from non-synchronous data'', \emph{Bernoulli} 19(2),
2013; N.~Huth and F.~Abergel, ``High frequency lead/lag relationships --- empirical facts'', \emph{Journal of Empirical Finance} 26, 2014.
\item Kenneth R. French Data Library, ``Portfolios formed on size [daily]''.
\end{omsources}

\section{Exercises}

\begin{exercise}[$\star$]\label{exo:rs:cross-sectional-and-cross-asset-features:1}
A group of five stocks returned 1\%, 2\%, 3\%, 4\% and 10\% today. Compute the fifth stock's return relative to its peers with and
without itself, and check \cref{prop:rs:cross-sectional-and-cross-asset-features:loo}.
\end{exercise}

\begin{exercise}[$\star$]\label{exo:rs:cross-sectional-and-cross-asset-features:2}
Among 486\,506 ordered pairs of which 133 are true links, how many true links does a random top 10\,000 contain on average? And a top
1\,000?
\end{exercise}

\begin{exercise}[$\star$]\label{exo:rs:cross-sectional-and-cross-asset-features:3}
A supplier's return includes 10\% of its customer's specific shocks, spread evenly over the next 21 days. A customer falls 6\% on
specific news today. What does the supplier's expected return drift by over the next 21 days, and per day?
\end{exercise}

\begin{exercise}[$\star\star$]\label{exo:rs:cross-sectional-and-cross-asset-features:4}
Two Brownian prices, $B(t) = A(t - L)$, are sampled on a grid of step $\Delta \ge L$. Show that the correlation of $B$'s return with
$A$'s previous return is $L/\Delta$, and that of their same-interval returns $1 - L/\Delta$.
\end{exercise}

\begin{exercise}[$\star\star$]\label{exo:rs:cross-sectional-and-cross-asset-features:5}
Why does the Hayashi--Yoshida estimator, rather than returns sampled on a grid, suit the estimation of a lead shorter than the average
time between trades?
\end{exercise}

\begin{exercise}[$\star\star$]\label{exo:rs:cross-sectional-and-cross-asset-features:6}
The small-stock quintile's weekly return correlates 0.31 with the large quintile's previous week and 0.37 with its own. Explain, with a
regression, why the lead may add little, and how the chapter tested it.
\end{exercise}

\begin{exercise}[$\star\star\star$]\label{exo:rs:cross-sectional-and-cross-asset-features:7}
\emph{Coding.} Keep the monthly customer-momentum feature but let the planted diffusion last 5 days, then 63 days
(\texttt{link\_days}). Compute the feature's rank IC and the truth's in each case. What does the result say about choosing a feature's
clock?
\end{exercise}

\begin{exercise}[$\star\star\star$]\label{exo:rs:cross-sectional-and-cross-asset-features:8}
\emph{Find the flaw.} ``The Nikkei's daily return is strongly correlated with the same day's S\&P 500 return, so we trade Japanese
stocks at the Tokyo close on the S\&P's move.''
\end{exercise}

\section{Problem: Who Moves First?}

\begin{problem}[{Weekend problem --- a planted lead, recovered and priced}]\label{pb:rs:cross-sectional-and-cross-asset-features:1}
Two instruments simulated by \texttt{firm.tape.simulate\_pair} on one efficient price, the second delayed by 0.5 seconds; then the
same at latencies 0, 0.25, 1 and 2 seconds and three seeds; then a slow version of the pair.

\medskip\noindent\textbf{Part I --- The cross-correlation.}
\begin{enumerate}
\item How often does each mid change in the fast pair, and in the slow one?
\item Why normalise the Hayashi--Yoshida cross-covariance by one-minute realised variances and not tick-level ones?
\item Where is the argmax of the fast pair's cross-correlation, and why is it not at 0.5?
\item What does the symmetric centre give, and on what assumption does it rest?
\item What is the lead--lag ratio of the fast pair at 0.5 seconds, and at zero latency?
\end{enumerate}

\medskip\noindent\textbf{Part II --- Recovery.}
\begin{enumerate}[resume]
\item What are the mean and worst errors of the centre and of the argmax on the fast pair?
\item Why do both fail on the slow pair?
\item What would you need to measure a 50-millisecond lead on a real pair?
\end{enumerate}

\medskip\noindent\textbf{Part III --- The feature.}
\begin{enumerate}[resume]
\item At a two-second interval, give the follower's correlation with the leader's last return, the reverse, and the follower's own.
\item Why are all three large on this simulator?
\item What is the lead's footprint at 2, 5, 10 and 30 seconds?
\item Why does the footprint vanish at thirty seconds, and why is it weak at 0.1 second?
\item How does this parallel the weekly size-quintile result?
\item Would the feature pay for crossing a one-tick spread? What does Huth and Abergel's evidence suggest?
\end{enumerate}

\medskip\noindent\textbf{Part IV --- Daily links.}
\begin{enumerate}[resume]
\item What are the customer-momentum IC, its $t$ statistic, and the IC of the truth?
\item Why is the feature's IC about half the truth's?
\item How many planted links are in the top 10\,000 mined pairs, daily and monthly, against the random expectation?
\item State the \emph{named result}: the estimated lead against the planted latency, and the share of cross-venue predictability left at
each sampling interval.
\item Which is easier to find without prior knowledge, a lead of half a second between two instruments on one price, or a one-month lead
between a customer and its supplier? Why?
\item In one sentence: what makes a \omterm{def:rs:cross-sectional-and-cross-asset-features:leadlag}{lead--lag relationship} a feature and not an artefact?
\end{enumerate}
\end{problem}

\section{Interview questions}

\begin{interviewq}[$\star$ \iqroles{researcher}]\label{iq:rs:cross-sectional-and-cross-asset-features:1}
How would you build a \omterm{def:rs:cross-sectional-and-cross-asset-features:peer}{peer group} for a stock, and what can go wrong with a \omterm{def:rs:cross-sectional-and-cross-asset-features:peer}{peer-relative return}?
\end{interviewq}

\begin{interviewq}[$\star\star$ \iqroles{researcher, trader}]\label{iq:rs:cross-sectional-and-cross-asset-features:2}
How would you estimate the lead of a future over an ETF from tick data?
\end{interviewq}

\begin{interviewq}[$\star\star$ \iqroles{researcher}]\label{iq:rs:cross-sectional-and-cross-asset-features:3}
Large stocks' returns predict small stocks' next-week returns. Is that a lead--lag effect?
\end{interviewq}

\begin{interviewq}[$\star\star$ \iqroles{researcher}]\label{iq:rs:cross-sectional-and-cross-asset-features:4}
You compute the lagged correlations of all pairs among 1\,000 stocks and trade the top 100. What do you expect?
\end{interviewq}

\begin{interviewq}[$\star\star$ \iqroles{researcher, trader}]\label{iq:rs:cross-sectional-and-cross-asset-features:5}
Why might a supplier's stock react slowly to its customer's news, and how would you test it?
\end{interviewq}

\begin{interviewq}[$\star\star\star$ \iqroles{researcher}]\label{iq:rs:cross-sectional-and-cross-asset-features:6}
Two prices, one a delayed copy of the other, are sampled at interval $\Delta$. What share of the follower's covariance with the leader
is predictable, as a function of $\Delta$ and the delay?
\end{interviewq}
